\section{Processi di supporto}
I processi di supporto includono l'insieme delle attività che sostengono e affiancano i processi primari e organizzativi, assicurando la corretta esecuzione delle operazioni, la tracciabilità delle scelte e la qualità degli artefatti lungo l'intero ciclo di vita del progetto.

\subsection{Documentazione}

\subsubsection{Introduzione}
Il processo di documentazione ha lo scopo primario di formalizzare, strutturare e mantenere tutte le informazioni prodotte dai processi del ciclo di vita del software adottati dal gruppo \textit{Git Happens}. In conformità con lo standard ISO/IEC 12207:1995, ogni attività svolta dal team deve essere adeguatamente pianificata, tracciata e registrata, al fine di garantire trasparenza, continuità operativa, manutenibilità e supporto al processo di valutazione della qualità.

\noindent La documentazione costituisce il principale mezzo di comunicazione formale del team, consentendo la piena comprensione delle decisioni e delle soluzioni adottate anche senza richiedere una costante mediazione orale, tracciando un chiaro confine tra creatività progettuale e disciplina metodologica. Gli obiettivi perseguiti da questo processo comprendono:
\begin{itemize}
    \item Definire procedure standardizzate e ripetibili per garantire uniformità strutturale e stilistica a tutta la documentazione di gruppo;
    \item Raccogliere e formalizzare norme operative condivise che agevolino e guidino i membri del team durante la stesura dei testi;
    \item Garantire la trasparenza e la verificabilità dei processi produttivi, attestando il raggiungimento della qualità attesa.
\end{itemize}
\noindent Il gruppo adotta la filosofia \textit{Documentation as Code}, trattando i documenti con il medesimo rigore riservato al codice sorgente. Tutti i documenti formali vengono redatti mediante il sistema di composizione tipografica \LaTeX{} e versionati all'interno del repository GitHub del team, assicurando rigore stilistico, controllo di versione e pubblicazione automatizzata. Il ciclo di vita di ciascun documento vede la cooperazione strutturata di tre ruoli distinti: il \textit{Redattore} per la stesura, il \textit{Verificatore} per il controllo di qualità e il \textit{Responsabile} per l'approvazione formale.

\subsubsection{Struttura dei documenti}
Tutti i documenti redatti dal gruppo seguono uno schema strutturale rigoroso e uniforme, volto a facilitare la consultazione. Ogni documento è composto da:
\begin{itemize}
    \item \textbf{Prima pagina}: contenente le informazioni identificative del gruppo, il titolo del documento, l'elenco dei redattori, dei verificatori e dei destinatari;
    \item \textbf{Registro delle modifiche}: una tabella di riepilogo storico posta prima dell'indice, recante versione, data, autore, verificatore e una descrizione sintetica della modifica;
    \item \textbf{Indice}: elenco gerarchico delle sezioni e sottosezioni, completo di indici delle figure ove applicabile;
    \item \textbf{Corpo del documento}: suddiviso in sezioni numerate progressivamente;
    \item \textbf{Pié di pagina}: contenente il nome del gruppo, il titolo del documento e la numerazione delle pagine.
\end{itemize}

\subsubsection{Convenzioni di nomenclatura}
Al fine di garantire uniformità, tracciabilità e una corretta automazione dei flussi di lavoro all'interno delle repository del gruppo \textbf{Git Happens}, vengono applicate rigorose convenzioni di nomenclatura per la gestione dei branch e dei file di documentazione.

\subsubsection{Nominazione dei branch di lavoro (Git)}
Il gruppo adotta il modello di sviluppo \emph{GitHub flow}. Per qualsiasi attività o contribuzione all'interno delle repository, il nome del branch deve essere generato seguendo una struttura standard legata al sistema di tracciamento dei ticket. 
La convenzione di naming ufficiale è la seguente:
\begin{center}
    \texttt{KEY-NUMBER-DESC}
\end{center}
dove:
\begin{itemize}
    \item \textbf{KEY}: rappresenta la chiave del progetto Jira di riferimento (ad esempio \texttt{SWE} per l'amministrazione e la documentazione del team);
    \item \textbf{NUMBER}: è il numero identificativo progressivo della issue su Jira;
    \item \textbf{DESC}: è una descrizione sintetica del task in formato minuscolo, con le parole separate da trattini alti (\texttt{-}).
\end{itemize}
\emph{Esempio}: \texttt{SWE-37-bootstrap-del-way-of-working}. 
L'utilizzo dell'estensione ufficiale di Atlassian per VS Code consente di automatizzare la creazione del branch e la corretta applicazione di questa convenzione direttamente dall'interfaccia di sviluppo.

\subsubsection{Nominazione dei file e struttura della documentazione}
La stesura e la pubblicazione della documentazione seguono uno standard organizzativo basato su cartelle dedicate all'interno della repository di progetto (ad esempio, la suddivisione temporale o tipologica per i verbali in \texttt{verbali/esterni/yyyy-mm-dd} o \texttt{verbali/interni/yyyy-mm-dd}).

Per quanto riguarda i file sorgente e gli asset correlati, valgono le seguenti regole:
\begin{itemize}
    \item \textbf{File principale}: ogni documento deve essere strutturato attorno a un file radice denominato \texttt{main.tex};
    \item \textbf{Capitoli e sezioni}: i file relativi ai singoli capitoli devono essere posizionati all'interno della cartella \texttt{capitoli/} e denominati con un prefisso numerico progressivo seguito da una descrizione chiara in minuscolo (es. \texttt{capitoli/01\_introduzione.tex});
    \item \textbf{Risorse multimediali}: le immagini e gli elementi grafici di supporto devono essere raccolti all'interno di una cartella denominata \texttt{assets} situata allo stesso livello del file \texttt{main.tex};
    \item \textbf{Metadati di pubblicazione}: ove previsto per migliorare la visualizzazione all'interno del sito web di documentazione del team, deve essere affiancato un file descrittivo in formato JSON denominato \texttt{metadata.json}.
\end{itemize}

\subsubsection{Versionamento}
Il controllo di versione costituisce un caposaldo della gestione della configurazione per il gruppo \textbf{Git Happens}, sia per quanto riguarda la documentazione sia per il codice sorgente del progetto. 
Il versionamento dei documenti segue la tripla numerica \texttt{X.Y.Z}, dove:
\begin{itemize}
    \item \textbf{X} viene incrementato in corrispondenza di rilasci fondamentali o baseline ufficiali concordate con i committenti e il proponente (es. RTB, PB);
    \item \textbf{Y} viene incrementato in caso di modifiche sostanziali, aggiunta di sezioni rilevanti o riscritture importanti del contenuto;
    \item \textbf{Z} viene incrementato in presenza di correzioni minori, rifiniture formali o aggiornamenti marginali.
\end{itemize}

Per quanto concerne il codice sorgente e la contribuzione alle repository, il team adotta la convenzione dei \emph{semantic commit messages} (messaggi di commit semantici scritti in lingua inglese). I messaggi devono esplicitare chiaramente la natura dell'operazione attraverso prefissi standardizzati quali:
\begin{itemize}
    \item \texttt{feat}: sviluppo o implementazione di una nuova funzionalità del progetto;
    \item \texttt{docs}: creazione o aggiornamento della documentazione;
    \item \texttt{chore}: attività di manutenzione, installazione di pacchetti, pulizia o organizzazione delle cartelle;
    \item \texttt{refactor}: modifiche strutturali del codice che non alterano le funzionalità esterne;
    \item \texttt{test}: scrittura, revisione o ampliamento dei test automatici (unitari e di integrazione).
\end{itemize}

\subsubsection{Procedure di redazione}
Il ciclo di vita di un documento o di un'attività all'interno del gruppo segue un iter rigoroso e tracciato, incentrato sull'uso di Atlassian Jira come unica fonte di verità per lo stato del progetto e sull'integrazione con GitHub:
\begin{enumerate}
    \item \textbf{Pianificazione (Jira)}: le attività vengono censite e gestite all'interno dello spazio Jira del team (identificato dalla chiave \texttt{SWE} per la documentazione o dallo spazio specifico del capitolato).
    \item \textbf{Apertura del branch}: a partire dal branch principale (\texttt{main}), lo sviluppatore crea un nuovo ramo di lavoro denominato secondo lo standard \texttt{KEY-NUMBER-DESC} (sfruttando eventualmente l'estensione Atlassian per VS Code per automatizzare l'operazione e il passaggio dello stato del ticket a \emph{In corso}).
    \item \textbf{Redazione e Commit}: si procede allo sviluppo o alla stesura del documento (strutturando i file in LaTeX tramite il file radice \texttt{main.tex} e inserendo i capitoli nella cartella dedicata), applicando la regola dei commit semantici.
    \item \textbf{Revisione e Pull Request}: completata la stesura, il ticket viene spostato in \emph{In revisione} su Jira, e viene aperta una \emph{Pull Request} (PR) su GitHub richiedendo la revisione dei verificatori assegnati.
    \item \textbf{Verifica e Build automatica}: i revisori esaminano la PR. Un flusso di Continuous Integration basato su GitHub Actions si attiva automaticamente per verificare la correttezza della build LaTeX.
    \item \textbf{Merge e Pubblicazione}: a seguito dell'approvazione formale e del superamento dei controlli automatici, viene effettuato il merge della PR nel ramo principale (spostando il ticket Jira in \emph{Completato}). Un ulteriore workflow di automazione provvede alla pubblicazione finale del documento all'interno del sito web di documentazione del gruppo.
\end{enumerate}

\subsection{Verifica}
Il processo di verifica garantisce che i prodotti sviluppati (documentazione e codice) siano conformi agli standard di qualità stabiliti. Esso si avvale di:
\begin{itemize}
    \item \textbf{Analisi statica}: controllo formale e strutturale eseguito mediante linters, type-check e l'ausilio di flussi automatici in CI per prevenire errori sintattici o violazioni delle norme di codifica.
    \item \textbf{Analisi dinamica e Testing}: esecuzione di test automatici (unitari, di integrazione e di sistema) per validare il comportamento effettivo del software.
    \item \textbf{Strumenti a supporto}: integrazione continua tramite GitHub Actions e utilizzo di strumenti di preview e controllo integrati nell'ambiente di sviluppo (es. estensione LaTeX Workshop per VS Code).
\end{itemize}

\subsection{Validazione}
La validazione consiste nella dimostrazione oggettiva che il prodotto finale risponda alle esigenze reali del committente e dell'azienda proponente Halue S.r.l. SB. Tale processo si realizza attraverso sessioni di riscontro periodico (come i SAL bisettimanali) e l'esecuzione di test di accettazione mirati a verificare il soddisfacimento dei requisiti funzionali e di qualità definiti nel Piano di Qualifica.

\subsection{Gestione della configurazione}
\subsubsection{Introduzione}
La gestione della configurazione disciplina il controllo e la tracciabilità di ogni elemento di configurazione (Configuration Item) del progetto. 
Le attività di tracciamento operativo si basano sulla sinergia tra l'Issue Tracking System di Jira (fonte di verità per la pianificazione e lo stato dei task) e GitHub (per il controllo di versione del codice e dei documenti). 

La gestione dei flussi di lavoro adotta il modello \emph{GitHub flow}, garantendo che ogni modifica sia isolata in appositi branch, sottoposta a revisione paritaria tramite Pull Request e validata da controlli automatici prima dell'incorporazione nel ramo stabile. 
Il coordinamento interno del team \textbf{Git Happens} e le comunicazioni esterne con il proponente completano l'infrastruttura di governo del progetto, i cui indicatori di qualità e le metriche di riferimento trovano riscontro puntuale all'interno del documento \textbf{Piano di Qualifica}.